\section*{Bab \ref{ch:b3:complex-mechanisms} --- Mekanisme Reaksi Kompleks}

\begin{solution}{exo:b3:complex-mechanisms:1}
Inert: $\ce{[Cr(H2O)6]^{3+}}$ ($d^3$) dan $\ce{[Co(NH3)6]^{3+}}$ ($d^6$ spin rendah), dengan energi aktivasi medan ligan
yang besar. \omterm{def:b3:complex-mechanisms:labile}{Labil}: $\ce{[Cr(H2O)6]^{2+}}$ ($d^4$ spin tinggi) dan $\ce{[Cu(H2O)6]^{2+}}$ ($d^9$),
yang terdistorsi Jahn--Teller dengan ikatan aksial yang panjang dan lemah; $\ce{[Zn(H2O)6]^{2+}}$ ($d^{10}$), tanpa penghalang
medan ligan.
\end{solution}

\begin{solution}{exo:b3:complex-mechanisms:2}
Jika $k_2[\mathrm Y] \gg k_{-1}[\mathrm X]$, penyebutnya adalah $k_2[\mathrm Y]$ dan $v = k_1[\mathrm{ML_5X}]$; jika $k_{-1}[\mathrm X] \gg k_2[\mathrm Y]$, penyebutnya adalah
$k_{-1}[\mathrm X]$ dan bentuk kedua mengikutinya (prakesetimbangan yang diinhibisi oleh ligan
yang pergi).
\end{solution}

\begin{solution}{exo:b3:complex-mechanisms:3}
D: $\Delta V^{\ddagger} > 0$ dan $\Delta^{\ddagger}S^\circ > 0$ (sebuah ligan lepas). A: keduanya negatif (sebuah ligan terikat dalam
keadaan transisi).
\end{solution}

\begin{solution}{exo:b3:complex-mechanisms:4}
$cis$-$\ce{[PtCl2(NH3)2]}$ dari $\ce{[PtCl4]^{2-}}$; $trans$-$\ce{[PtCl2(NH3)2]}$ dari $\ce{[Pt(NH3)4]^{2+}}$ (\omterm{def:b3:complex-mechanisms:trans-effect}{efek
trans} \ce{Cl-} melampaui \omterm{def:b3:complex-mechanisms:trans-effect}{efek trans} \ce{NH3}).
\end{solution}

\begin{solution}{exo:b3:complex-mechanisms:5}
$k_{\mathrm{obs}} = 2.0 \times \num{3.0e4} \times 0.10/1.2 = \qty{5.0e3}{s^{-1}}$; pada \qty{1.0}{mol/L}, $\num{6.0e4}/3.0 = \qty{2.0e4}{s^{-1}}$; limitnya
$k_i = \qty{3.0e4}{s^{-1}}$.
\end{solution}

\begin{solution}{exo:b3:complex-mechanisms:6}
$\Delta V^{\ddagger} = -RT\ln2/\Delta p = -8.314 \times 298 \times 0.693/\num{99.9e6} = \qty{-1.7e-5}{m^3.mol^{-1}} = \qty{-17}{cm^3.mol^{-1}}$:
asosiatif.
\end{solution}

\begin{solution}{exo:b3:complex-mechanisms:7}
Fosfina, donor $\sigma$ dan akseptor $\pi$ yang kuat, melemahkan ikatan yang trans terhadapnya (\omterm{def:b3:complex-mechanisms:trans-effect}{pengaruh
trans}, ikatan lebih panjang) dan melabilkannya (\omterm{def:b3:complex-mechanisms:trans-effect}{efek trans}): klorida itulah yang
digantikan lebih dahulu.
\end{solution}

\begin{solution}{exo:b3:complex-mechanisms:8}
Ligan penjembatan akhirnya berpindah ke logam yang teroksidasi; laju sangat
bergantung pada jenis calon jembatan; sebuah zat antara dwiinti terdeteksi; dan
laju tidak dapat melampaui laju substitusi yang diperlukan untuk membentuk
jembatan pada pasangan yang \omterm{def:b3:complex-mechanisms:labile}{labil}.
\end{solution}

\begin{solution}{exo:b3:complex-mechanisms:9}
$(\lambda + \Delta G^\circ)^2/4\lambda$ dengan $4\lambda = 320$: 20.0, 11.3, 0, dan \qty{5.0}{kJ/mol} (yang terakhir berada di \omterm{def:b3:complex-mechanisms:inverted}{daerah
terbalik}).
\end{solution}

\begin{solution}{exo:b3:complex-mechanisms:10}
$k_{12} = \sqrt{4.0 \times \num{2.0e5} \times \num{1.0e4}} = \qty{8.9e4}{L.mol^{-1}.s^{-1}}$.
\end{solution}

\begin{solution}{exo:b3:complex-mechanisms:11}
$d^5$ spin tinggi: 0; $d^7$: $1.33\,Dq$; $d^8$: $2.0\,Dq$. Laju pertukaran air yang diharapkan: \ce{Mn^{2+}} $>$ \ce{Co^{2+}} $>$
\ce{Ni^{2+}}, yaitu urutan yang teramati.
\end{solution}

\begin{solution}{exo:b3:complex-mechanisms:12}
Untuk reaksi yang sangat eksergonik, laju intrinsiknya melampaui laju pertemuan
ion-ion: tetapan yang teramati tetap pada batas difusi dan penurunannya
tersembunyi. Dengan donor dan akseptor yang dipertahankan pada jarak tertentu
dalam satu molekul, tidak ada tahap difusi, dan serangkaian akseptor dengan
gaya dorong yang makin besar memperlihatkan laju yang naik, melewati maksimum,
lalu turun.
\end{solution}

\begin{solution}{pb:b3:complex-mechanisms:1}
\textbf{1.} $5d^8$.
\textbf{2.} Medan yang kuat pada ion $5d$ membuat orbital $d_{x^2-y^2}$, yang mengarah ke empat
ligan, sangat tinggi: delapan elektron berpasangan di empat orbital yang lebih rendah.
\textbf{3.} Kompleks itu memiliki 16 elektron valensi dan sebuah situs aksial yang terbuka: ligan yang
masuk terikat lebih dahulu, melalui keadaan transisi berkoordinasi lima.
\textbf{4.} $k_{\mathrm{obs}} = k_1 + k_2[\mathrm Y]$: jalur pelarut (serangan pelarut, lalu penggantian cepat oleh Y)
dan jalur langsung.
\textbf{5.} Substitusinya memakan waktu berjam-jam (Bagian III): inert dalam skala menit, tetapi tidak dalam
skala satu hari.
\textbf{6.} Cisplatin, karena \ce{NH3} kedua menggantikan klorida yang trans terhadap klorida
(\omterm{def:b3:complex-mechanisms:trans-effect}{efek trans} $\ce{Cl-} > \ce{NH3}$); produk samping juga terbentuk.
\textbf{7.} Iodida memiliki \omterm{def:b3:complex-mechanisms:trans-effect}{efek trans} yang lebih besar daripada klorida, sehingga efek pengarahnya
lebih bersih dan lebih cepat.
\textbf{8.} Dalam $\ce{[PtI3(NH3)]-}$, iodida yang trans terhadap iodida-lah yang dilabilkan, bukan yang trans terhadap
\ce{NH3}: amonia kedua masuk pada posisi cis terhadap amonia pertama.
\textbf{9.} Perak iodida mengendap dan $cis$-$\ce{[Pt(H2O)2(NH3)2]^{2+}}$ tetap berada dalam larutan.
\textbf{10.} Klorida menggantikan kedua molekul air pada posisinya masing-masing; ligan-ligan
amin tidak tersentuh.
\textbf{11.} Dari $\ce{[Pt(NH3)4]^{2+}}$ dan dua klorida.
\textbf{12.} Kedua posisi reaktifnya saling berseberangan: transplatin tidak dapat mengikat dua
basa bertetangga pada untai DNA yang sama, yaitu lesi yang menjadi jalan kerja
cisplatin.
\textbf{13.} $\ln2/\num{9.0e-5} = \qty{7.7e3}{s} = \qty{2.1}{h}$.
\textbf{14.} $1 - \eu^{-\num{9.0e-5} \times 7200} = 0.48$.
\textbf{15.} Air adalah gugus pergi yang jauh lebih baik daripada klorida: kompleks akua
dengan cepat digantikan oleh nitrogen sebuah basa DNA.
\textbf{16.} Asosiatif, seperti semua substitusi platina(II): \omterm{def:b3:complex-mechanisms:activation-volume}{volume aktivasinya}
negatif.
\textbf{17.} Konsentrasi klorida yang tinggi mendorong kesetimbangan kembali ke arah
bentuk dikloro, sehingga obat itu tetap utuh sampai mencapai sel.
\textbf{18.} $f = K/(K + [\ce{Cl-}])$.
\textbf{19.} $\num{3.0e-3}/0.103 = 0.029$.
\textbf{20.} $\num{3.0e-3}/0.0070 = 0.43$.
\textbf{21.} 15.
\textbf{22.} Hanya di dalam sel terdapat fraksi yang berarti dalam bentuk kompleks akua
yang reaktif.
\textbf{23.} Ligan-ligan belerang, seperti glutation serta sisteina dan metionina
protein, yang mengikat platina yang lunak dengan kuat.
\textbf{24.} Amonia terikat kuat pada platina dan terletak trans terhadap klorida, ligan yang
\omterm{def:b3:complex-mechanisms:trans-effect}{efek trans-nya} kecil: ikatan-ikatan amin tidak dilabilkan.
\textbf{25.} \textbf{Fraksi obat dalam bentuk akua yang reaktif sekitar 15 kali lebih besar di dalam sel
daripada di dalam plasma.}
\end{solution}
